\chapter{数值实验与结果分析}
\section{环境、指标与数据来源}
实际运行环境为Linux CPU、Python 3.12.14、NumPy 2.2.6、SciPy 1.15.3及Matplotlib 3.10.3，BLAS线程设为1。程序设计为Windows Python 3.12可运行，但尚未完成Windows实机验证。无需GPU或真实外部数据。直线模型多数使用五种子，花形近似复现主比较使用十种子；表中的均值覆盖全部指定运行，样本标准差采用分母$n-1$。

\begin{table}[htbp]\centering\small
\caption{主要实验设置概览，细项以各run配置为准}\label{tab:protocol_overview}
\begin{tabular}{lp{.69\textwidth}}\toprule
设置 & 固定参数与比较范围\\\midrule
平滑采样主比较 & 每侧64特征，scale2；三轮；A/B/C/D及两消融；三系数组合；点预算1104或行预算1232；五种子。\\
局部峰比较 & 每侧128特征，scale2/4；峰宽.12/.06；三轮/六轮；1232/1904行；两个系数组合；五种子。\\
花形主近似复现 & 每侧80/160/320特征；总配置点1000/1000/5000；两网络权重范围1.6、偏置.7；FD和解析导数；十种子。\\
\bottomrule\end{tabular}\end{table}

直线模型测试每侧83乘83中点及257界面点；峰值细测试每侧161乘161、界面521点，另用81乘81检查积分敏感性。全局Poisson随机特征使用81乘81内部点。FDM误差在其节点上，残差为离散系统相对残差；其他PDE残差为独立点上的连续解析算子残差，两者不能直接作为同指标比较。时间区分单次组装/求解和完整自适应过程，不含最终验证/绘图。

所有表格由原始CSV生成，图由CSV重绘。主要来源为实验01、02、03、04、06、07、08及09的两组复现记录，均为2026年10月13日保存。完整路径见附录；原文数据另标注，不混为本项目输出。

设同一设置下五次误差为$e_s$，报告均值与样本标准差：
\[\bar e=\frac1n\sum_s e_s,\qquad s_e=\sqrt{\frac{\sum_s(e_s-\bar e)^2}{n-1}}.\]
标准差描述随机初始化的离散程度，不等于均值的置信区间。采样比较采用共享种子的配对比例$R_s=e_{D,s}/e_{X,s}$，其中$X$为基线；$R_s<1$仅表示该次解误差更低。比例中位数与均值之比不同，本文不混用。五种子不足以支持广泛的统计显著性或总体优势声明。

\section{Poisson基线正确性}
表\ref{tab:fdm}和图\ref{fig:fdm}显示四级网格的误差与观测阶。相邻步长减半后误差约降四倍，与第二章的正弦特征值推导一致。单元测试还核查矩阵符号、制造右端、tanh导数和秩亏最小范数解。
\begin{table}[htbp]\centering\caption{Poisson五点差分真实结果}\label{tab:fdm}\begin{tabular}{lll}\toprule
每向内部点 & 相对离散误差 & 观测阶 \\ \midrule
15 & $3.219\times10^{-3}$ & -- \\
31 & $8.036\times10^{-4}$ & 2.0021 \\
63 & $2.008\times10^{-4}$ & 2.0005 \\
127 & $5.020\times10^{-5}$ & 2.0001 \\
\bottomrule\end{tabular}
\end{table}
\figurefile{fdm_convergence}{正弦Poisson五点差分收敛；数据来自exp01原始CSV。}{fig:fdm}
基础全局RaNN每种特征数五次：64特征相对误差均值约$8.64\times10^{-3}$，128特征约$5.99\times10^{-5}$，后者样本标准差约$3.08\times10^{-5}$。与此同时128特征加权矩阵平均条件数约$2.48\times10^9$。精度改善并不意味着条件性改善。全局网络与双子域网络总特征数、坐标尺度、配置点及评估方式不同，本文不据此宣称公平效率优越。

\section{直线界面与尺度取舍}
150次运行包含连续与非零跳跃、七组正反向系数组合、raw/balanced及Poisson极限。无运行异常，平衡组通过预设粗误差门槛；这一门槛不覆盖后续窄峰。表\ref{tab:interface}为非零跳跃、每侧128特征、scale2的五次均值。
\begin{table}[htbp]\centering\caption{高对比度下解与物理通量，五种子均值及样本标准差}\label{tab:interface}\resizebox{\textwidth}{!}{\begin{tabular}{llll}\toprule
$(\beta_1,\beta_2)$ & 尺度 & 相对$L^2$：均值$\pm s$ & 通量RMS：均值$\pm s$ \\ \midrule
(1000,1) & raw & $(4.579\pm4.578)\times10^{-5}$ & $(1.763\pm2.570)\times10^{-4}$ \\
(1000,1) & balanced & $(1.391\pm0.681)\times10^{-7}$ & $(2.756\pm3.733)\times10^{-4}$ \\
(1,1000) & raw & $(1.094\pm0.898)\times10^{-4}$ & $(1.711\pm1.554)\times10^{-4}$ \\
(1,1000) & balanced & $(1.466\pm0.280)\times10^{-7}$ & $(4.453\pm2.523)\times10^{-4}$ \\
\bottomrule\end{tabular}
}\end{table}
\figurefile{interface_contrast}{非零跳跃、每侧128特征的系数对比度实验；物理通量未除以系数。}{fig:interface}
平衡后解误差明显降低，但物理通量均值不一定降低。raw最高条件数约$1.69\times10^{15}$、最低有效秩240/256；balanced仍可达到约$2.11\times10^{12}$、最低255/256。截断SVD使近秩亏系统可求解，但不能据正常运行声称病态已消除。

\section{采样主比较与消融}
180次主比较固定每侧64特征、scale2、三轮；每种预算覆盖三个系数组合和五种子。表\ref{tab:sampling}展示点预算、$\beta=(1000,1)$的全部六种方法，最终物理总点1104，D方程数可变。
\begin{table}[htbp]\centering\caption{等物理点预算下采样与消融，误差为均值及样本标准差}\label{tab:sampling}\resizebox{\textwidth}{!}{\begin{tabular}{llll}\toprule
方法 & 相对$L^2$：均值$\pm s$ & 完整时间/秒 & 界面点均值 \\ \midrule
A & $(1.271\pm0.411)\times10^{-4}$ & 0.0133 & 128.0 \\
B & $(1.200\pm0.404)\times10^{-4}$ & 0.0137 & 128.0 \\
C & $(1.555\pm0.542)\times10^{-4}$ & 0.0763 & 128.0 \\
D & $(1.591\pm0.541)\times10^{-4}$ & 0.0767 & 68.0 \\
去PDE指标 & $(1.103\pm0.253)\times10^{-4}$ & 0.0829 & 512.0 \\
去界面指标 & $(1.561\pm0.545)\times10^{-4}$ & 0.0743 & 56.0 \\
\bottomrule\end{tabular}
}\end{table}
\figurefile{sampling}{点预算下全部种子的解误差，横线为均值；数据来自exp03。}{fig:sampling}
在每种预算的3设置乘5种子配对中，D对A/B解误差胜出均为0/15，D/B误差比例中位数约1.264和1.256。设置之间相关，不能把15个配对当作15个独立种子做显著性推断。D对C仅部分改善，仍没有稳定胜出。去PDE指标消融给界面更多点而较好，但后续相同最终组配额的固定随机在六组均值中更好且更快，未支持把收益简单归因于残差选点。

这说明比较基线必须合理：若只选界面明显欠采样的静态比例，可能高估动态分配的价值。另一方面，D默认指标偏向内部，界面配额较低，也可能影响解误差。本研究把该解释保留为数值猜想，而非因果证明。

表\ref{tab:paired_sampling}按系数与预算拆开配对结果，避免只报告合并中位数而掩盖不同设置。三个系数组合均使用相同五个种子，跨设置重复并不产生新的独立初始化样本。
\begin{table}[p]\centering\caption{平滑模型的分设置配对；比值为D误差除以基线误差}\label{tab:paired_sampling}\begin{tabular}{lllll}\toprule
预算 & $(\beta_1,\beta_2)$ & 配对 & 误差较小次数 & 比值中位数 \\ \midrule
点 & (1,1) & D/A & 0/5 & 1.2605 \\
点 & (1,1) & D/B & 0/5 & 1.2619 \\
点 & (1,1) & D/C & 0/5 & 1.0175 \\
点 & (1000,1) & D/A & 0/5 & 1.2338 \\
点 & (1000,1) & D/B & 0/5 & 1.2924 \\
点 & (1000,1) & D/C & 1/5 & 1.0308 \\
点 & (1,1000) & D/A & 0/5 & 1.2721 \\
点 & (1,1000) & D/B & 0/5 & 1.2130 \\
点 & (1,1000) & D/C & 3/5 & 0.9975 \\
行 & (1,1) & D/A & 0/5 & 1.2248 \\
行 & (1,1) & D/B & 0/5 & 1.2451 \\
行 & (1,1) & D/C & 2/5 & 1.0156 \\
行 & (1000,1) & D/A & 0/5 & 1.2339 \\
行 & (1000,1) & D/B & 0/5 & 1.3013 \\
行 & (1000,1) & D/C & 1/5 & 1.0183 \\
行 & (1,1000) & D/A & 0/5 & 1.2726 \\
行 & (1,1000) & D/B & 0/5 & 1.2561 \\
行 & (1,1000) & D/C & 2/5 & 1.0197 \\
\bottomrule\end{tabular}
\end{table}

\subsection{配额与选点位置的分离诊断}
去PDE指标消融在点预算下最终每侧内部176点、界面512点，外边界240点，共1104个物理点及1616个标量方程。后续固定随机对照采用同一最终配额、特征数及随机种子，表\ref{tab:allocation_control}显示其平均解误差在三个系数组合下均更低。该比较是看到消融结果后登记的探索性诊断，不能改写成预先设计的验证。表中成本来自不同批次，支持当前环境下的成本观察；严格串行重测的证据见峰值场景。
\begin{table}[htbp]\centering\caption{相同最终配额的随机对照与去PDE指标消融，五种子}\label{tab:allocation_control}\resizebox{\textwidth}{!}{\begin{tabular}{llll}\toprule
$(\beta_1,\beta_2)$ & 同配额方法 & 相对$L^2$：均值$\pm s$ & 完整时间/秒 \\ \midrule
(1,1) & 固定随机 & $(1.112\pm0.528)\times10^{-4}$ & 0.0168 \\
(1,1) & 去PDE指标 & $(1.156\pm0.594)\times10^{-4}$ & 0.0826 \\
(1000,1) & 固定随机 & $(1.074\pm0.299)\times10^{-4}$ & 0.0169 \\
(1000,1) & 去PDE指标 & $(1.103\pm0.253)\times10^{-4}$ & 0.0829 \\
(1,1000) & 固定随机 & $(9.311\pm5.096)\times10^{-5}$ & 0.0167 \\
(1,1000) & 去PDE指标 & $(9.922\pm5.283)\times10^{-5}$ & 0.0826 \\
\bottomrule\end{tabular}
}\end{table}
这个对照未支持残差排序选点对消融收益的必要性，但也不是完整因果分解：增量历史、位置与求解路径仍不同。因各组按点数取平均，增加界面点没有直接增加界面平方权重，而是改变界面经验测度并减少内部预算。不能将此同配额实验与1232行默认方案直接称为等行预算比较。

\section{参数、秩与误差饱和}
285次探索性后续实验包含固定比例及单因素扫描。表\ref{tab:width}为$\beta=(1,1000)$宽度对照，其他设置固定。
\begin{table}[htbp]\centering\caption{特征宽度敏感性，数据秩上限为总特征数}\label{tab:width}\begin{tabular}{llll}\toprule
每侧特征 & 相对$L^2$均值 & 原始条件数均值 & 数据秩范围 \\ \midrule
32 & $1.138\times10^{-2}$ & $1.823\times10^{5}$ & 64--64 \\
64 & $1.085\times10^{-4}$ & $1.516\times10^{8}$ & 128--128 \\
128 & $1.236\times10^{-7}$ & $9.579\times10^{11}$ & 255--256 \\
256 & $2.581\times10^{-10}$ & $1.076\times10^{17}$ & 404--417 \\
\bottomrule\end{tabular}
\end{table}
\figurefile{width}{三个系数组合下的宽度、物理通量与条件数；数据来自exp04。}{fig:width}
当前平滑模型增宽改善误差，但条件数升至约$10^{17}$，256/侧时有效秩仅404至417，不能以总特征数代表独立维度。scale从.5到8的影响显著，小范围可适合平滑函数却加重病态，大范围可能使近似质量变差。界面权重、岭和点数也改变排序：强岭增加偏差；点数216至1728的解误差不单调，物理通量却可改善。这些扫描属于看到主比较后登记的探索性研究，不把其中最好参数改写为事先最优基线。

\section{独立局部峰与多预算}
在原分片制造解两侧加$h=16A x(1-x)y(1-y)\exp(-((x-c_x)^2+(y-c_y)^2)/w^2)$，取$A=.75,(c_x,c_y)=(.46,.55)$，两峰宽$.12,.06$。值跳跃不变，但通量增加$(\beta_2-\beta_1)h_x$，须重新制造数据。400次比较采用每侧128特征、scale2/4、两个行预算及五方法五种子。两个测试网格误差相对变化最大约0.49\%，未进入选点。
\figurefile{peaks}{scale4下独立峰的两个标量行预算比较；所有预设模型与方法均保留。}{fig:peaks}
80描述性配对中，D对A、B、C误差更低分别6、36、49次；对应比值中位数约1.293、1.095、.963。D\_unused没有稳定逆转结果。部分窄峰相对误差大于1，最大约1.493，无运行异常并不代表精度验收。固定特征表达力和采样目标变化都可能影响结果，不能只增加点数就承诺解决。

主峰值前段与谱任务并发，严格成本讨论使用另行25次串行重测。表\ref{tab:peaks}为固定窄峰、$\beta=(1000,1)$、scale4、六轮及1904方程，系数已与主case校验一致。D在该设置比C精度更好，但仍不如A且更贵。
\begin{table}[htbp]\centering\caption{固定峰值场景的串行完整成本，误差均值及样本标准差}\label{tab:peaks}\resizebox{\textwidth}{!}{\begin{tabular}{lll}\toprule
方法 & 相对$L^2$：均值$\pm s$ & 串行完整时间/秒 \\ \midrule
A & $(1.948\pm0.610)\times10^{-1}$ & 0.0481 \\
B & $(3.839\pm1.098)\times10^{-1}$ & 0.0450 \\
C & $(3.524\pm1.360)\times10^{-1}$ & 0.4037 \\
D & $(2.845\pm0.859)\times10^{-1}$ & 0.4097 \\
D\_unused & $(2.742\pm0.801)\times10^{-1}$ & 0.4122 \\
\bottomrule\end{tabular}
}\end{table}

窄峰结果并未唯一识别误差来源。表达空间、选点分布、组权重及数值截断均可能限制精度；尚未计算该空间的最佳近似误差，因此不能把失败严格归因于表达力。可以确定的是：在本文已运行的固定特征和预算下，联合采样没有实现稳定的精度或成本优势。

\section{固定矩阵扰动观察}
15组矩阵谱、三阈值和两个扰动方向构成90组观察。右端相对扰动固定$10^{-8}$；最小保留左奇异方向的放大与$1/\sigma_{\min,\rm retained}$在浮点误差内一致，图\ref{fig:perturb}展示系数和函数变化的差别。
\figurefile{perturbation}{固定TSVD空间下的右端扰动；理论界仅针对固定矩阵和保留空间。}{fig:perturb}
最大系数相对变化约275.6，而预测函数相对变化最大约$1.71\times10^{-8}$。冗余特征可使系数敏感而函数变化较小，但该数值还可能接近大系数抵消的浮点求值地板，未以高精度算术复核。不能由此断言病态无害，也没有考察矩阵扰动、阈值跨秩或所有物理数据误差。

\section{花形界面FD近似复现}
选Li--Wang预印本Example4.1\cite{interfacepreprint}，$\Omega=(-1,1)^2$，中心$.02\sqrt5$，极径$r=.4+.2\sin20\theta$。图1核查花内为$\Omega_1$。两侧真解为$e^{xy}/\beta_1$及$\sin x\sin y/\beta_2$，$\beta=(1,10)$；法向由曲线切向旋转并归一化，不能只用径向单位向量。

按最贴合文字的解释，两网络权重范围均为1.6、偏置范围.7。与原文相同的设置包括单隐层tanh、3:1:1点数比例、$h_1=10^{-6},h_2=5\times10^{-4}$、行权重50、GL20及十次平均。作者种子、dtype、最小二乘阈值、确切界面测度未获实现；本项目float64/gelsd/$\tau=10^{-12}$、主角度均匀并补弧长对照。因此\textbf{只称近似复现}。
\begin{table}[htbp]\centering\small\caption{原文表1与本项目近似复现，十种子均值}\label{tab:reproduction}\resizebox{\textwidth}{!}{\begin{tabular}{lllll}\toprule
$m/N$ & 导数 & 原文表1 & GL20：均值$\pm s$ & GL100：均值$\pm s$ \\ \midrule
80/1000 & fd & $1.180\times10^{-3}$ & $(3.169\pm2.090)\times10^{-5}$ & $(3.236\pm2.147)\times10^{-5}$ \\
80/1000 & analytic & $1.180\times10^{-3}$ & $(3.169\pm2.090)\times10^{-5}$ & $(3.236\pm2.147)\times10^{-5}$ \\
160/1000 & fd & $9.200\times10^{-6}$ & $(1.059\pm0.664)\times10^{-7}$ & $(1.080\pm0.683)\times10^{-7}$ \\
160/1000 & analytic & $9.200\times10^{-6}$ & $(1.059\pm0.664)\times10^{-7}$ & $(1.080\pm0.683)\times10^{-7}$ \\
320/5000 & fd & $3.170\times10^{-8}$ & $(3.057\pm0.712)\times10^{-10}$ & $(3.076\pm0.457)\times10^{-10}$ \\
320/5000 & analytic & $3.170\times10^{-8}$ & $(2.950\pm0.721)\times10^{-10}$ & $(2.961\pm0.439)\times10^{-10}$ \\
\bottomrule\end{tabular}
}\end{table}
\figurefile{reproduction}{最贴合文字初始化下的FD与解析导数对照；原文曲线不是本项目数据。}{fig:reproduction}
表\ref{tab:reproduction}的误差与原文明显不同。更低数值不能当作算法改进，差异可能涉及初始化、dtype、截断与采样，尚无法唯一归因。最初按子域解释范围的110次保留为独立对照，核对文字后追加60次权重/偏置解释，旧数据不改写。原文通量印刷矩阵疑似漏$\beta$，本项目按物理模型保留，正式版仍待核对。

在子域范围解释变体中，五种子的FD步长结果见表\ref{tab:steps}。过大步长有截断偏差，过小步长出现舍入放大；不能把此变体当作作者结论。GL20与GL100在极小误差下可有明显相对变化，均保留而不选择较好的积分值。
\begin{table}[htbp]\centering\caption{FD步长独立诊断，$m=160,N=1000$}\label{tab:steps}\begin{tabular}{lll}\toprule
二阶差分步长 & 相对$L^2$均值 & 解析PDE残差RMS \\ \midrule
$1.000\times10^{-2}$ & $2.635\times10^{-8}$ & $4.659\times10^{-6}$ \\
$5.000\times10^{-4}$ & $5.976\times10^{-10}$ & $2.213\times10^{-7}$ \\
$1.000\times10^{-5}$ & $1.262\times10^{-7}$ & $1.768\times10^{-5}$ \\
$1.000\times10^{-6}$ & $8.669\times10^{-6}$ & $8.555\times10^{-4}$ \\
\bottomrule\end{tabular}
\end{table}

定义积分敏感度为$\Delta=|e_{20}-e_{100}|/e_{100}$。表\ref{tab:quadrature}使用最贴合原文字面初始化的FD组，逐种子计算后报告中位数与最大值。GL100是更密的数值验证规则，不是连续误差的严格真值；改变规则只用于诊断，没有进入训练或选择最好运行。误差极小时较大的相对差别也可能对应很小的绝对差别，故同时报告误差量级。
\begin{table}[htbp]\centering\caption{花形界面误差对积分规则的敏感度，十种子}\label{tab:quadrature}\resizebox{\textwidth}{!}{\begin{tabular}{llll}\toprule
$m/N$ & 积分差别中位数 & 积分差别最大值 & GL100误差均值 \\ \midrule
80/1000 & 2.22\% & 3.38\% & $3.236\times10^{-5}$ \\
160/1000 & 1.52\% & 4.83\% & $1.080\times10^{-7}$ \\
320/5000 & 7.58\% & 33.30\% & $3.076\times10^{-10}$ \\
\bottomrule\end{tabular}
}\end{table}

\section{实验结论的范围}
已有数据支持对采样、尺度、表达力和求解截断进行系统讨论，但不支持普适自适应优势。平滑、峰值及花形模型具有不同的近似难度；对一个算例表现好的范围不应直接迁移。还未做传统界面FDM/FEM或PINN同预算对照、完整多层增长、复杂物理或矩阵扰动。初稿中所有结论限于已登记模型与参数，后续应以独立模型和理论补充，而非对旧测试集无限调参。

\section{多指标证据与采样过程分析}
\subsection{相同标量行预算下的独立误差}
为避免只根据相对解误差评价算法，表\ref{tab:multimetric}补列1232行主实验的最大绝对误差、两侧归一化PDE残差、界面值RMS和物理通量RMS。全部为五种子的算术均值；原始逐次数据与标准差仍保存于同一CSV及汇总表。归一化PDE残差除以各自扩散系数，而最后一列通量未作归一化，读表时不可把这些量直接按数值大小比较。
\begin{table}[p]\centering
\caption{相同1232行预算下的独立多指标均值；每侧64特征、scale2、三轮}\label{tab:multimetric}
\resizebox{\textwidth}{!}{\begin{tabular}{lllllll}\toprule
$(\beta_1,\beta_2)$ & 方法 & $E_\infty$ & $R_1/\beta_1$ & $R_2/\beta_2$ & $J_u$ & $J_q$ \\ \midrule
(1,1) & A & $4.895\times10^{-4}$ & $2.989\times10^{-3}$ & $3.125\times10^{-3}$ & $1.124\times10^{-4}$ & $9.465\times10^{-5}$ \\
(1,1) & B & $4.449\times10^{-4}$ & $3.866\times10^{-3}$ & $4.382\times10^{-3}$ & $1.185\times10^{-4}$ & $1.179\times10^{-4}$ \\
(1,1) & C & $5.999\times10^{-4}$ & $3.227\times10^{-3}$ & $3.010\times10^{-3}$ & $1.519\times10^{-4}$ & $1.236\times10^{-4}$ \\
(1,1) & D & $6.116\times10^{-4}$ & $2.974\times10^{-3}$ & $2.988\times10^{-3}$ & $1.452\times10^{-4}$ & $1.074\times10^{-4}$ \\
(1000,1) & A & $4.537\times10^{-4}$ & $3.108\times10^{-3}$ & $2.656\times10^{-3}$ & $1.123\times10^{-4}$ & $1.029\times10^{-1}$ \\
(1000,1) & B & $3.964\times10^{-4}$ & $4.613\times10^{-3}$ & $3.906\times10^{-3}$ & $1.093\times10^{-4}$ & $1.047\times10^{-1}$ \\
(1000,1) & C & $5.385\times10^{-4}$ & $3.075\times10^{-3}$ & $2.574\times10^{-3}$ & $1.450\times10^{-4}$ & $1.376\times10^{-1}$ \\
(1000,1) & D & $5.506\times10^{-4}$ & $3.036\times10^{-3}$ & $2.502\times10^{-3}$ & $1.436\times10^{-4}$ & $1.219\times10^{-1}$ \\
(1,1000) & A & $4.082\times10^{-4}$ & $2.665\times10^{-3}$ & $2.913\times10^{-3}$ & $7.398\times10^{-5}$ & $8.728\times10^{-2}$ \\
(1,1000) & B & $3.720\times10^{-4}$ & $3.402\times10^{-3}$ & $3.868\times10^{-3}$ & $7.790\times10^{-5}$ & $9.671\times10^{-2}$ \\
(1,1000) & C & $4.923\times10^{-4}$ & $2.733\times10^{-3}$ & $2.961\times10^{-3}$ & $9.955\times10^{-5}$ & $1.135\times10^{-1}$ \\
(1,1000) & D & $5.153\times10^{-4}$ & $2.655\times10^{-3}$ & $2.869\times10^{-3}$ & $9.029\times10^{-5}$ & $1.095\times10^{-1}$ \\
\bottomrule\end{tabular}
}
\end{table}

在$\beta=(1000,1)$设置中，D的右侧归一化PDE残差均值约$2.50\times10^{-3}$，低于B的约$3.91\times10^{-3}$；但D的最大绝对误差约$5.51\times10^{-4}$，高于B的约$3.96\times10^{-4}$。这提供了一个具体反例：在独立点上内部残差降低，也不能不加条件地推出最大解误差降低。原因可能同时涉及界面、边界、截断及不同误差传播方向，不能只根据这两列唯一归因。

同一设置下，D的物理通量RMS约0.122，B约0.105，说明残差目标在内部改善时并未同时改善全部界面指标。这个量级来自每侧64特征的采样主比较，不能与每侧128特征的界面验证表混用。若不注明特征数，读者可能误以为高对比度通量结果前后矛盾；实际上两表属于不同近似能力和配置点实验。

对于单位系数，D两侧PDE残差也处于较低水平，但界面值误差和最大解误差没有同步胜过静态基线。因此本研究的负面结果不仅是某个高对比度下的偶然异常。另一方面，不能把当前三个平滑设置概括为联合采样在所有问题上无效；峰值实验已经显示D相对仅PDE采样在部分设置改善，这种差异仍需按问题和指标分别报告。

\subsection{条件数较低为何仍可能误差较高}
表\ref{tab:rankcost}显示相同行预算的条件数、有效秩、物理点及时间。这里总列数为128，所有列出的主比较有效秩均为128，原始条件数约为$10^8$量级。D在高对比度设置下的平均条件数略低于B，却仍没有更小的解误差。这说明条件数反映最坏线性敏感性，不能直接作为近似误差的单调代理。
\begin{table}[p]\centering
\caption{1232行主比较的矩阵与资源诊断，五种子；时间含完整自适应过程}\label{tab:rankcost}
\resizebox{\textwidth}{!}{\begin{tabular}{lllllll}\toprule
$(\beta_1,\beta_2)$ & 方法 & $\kappa$均值 & 有效秩 & 物理点均值 & 界面点均值 & 时间/秒 \\ \midrule
(1,1) & A & $1.812\times10^{8}$ & 128--128 & 1104.0 & 128.0 & 0.0135 \\
(1,1) & B & $1.766\times10^{8}$ & 128--128 & 1104.0 & 128.0 & 0.0137 \\
(1,1) & C & $1.618\times10^{8}$ & 128--128 & 1104.0 & 128.0 & 0.0752 \\
(1,1) & D & $1.636\times10^{8}$ & 128--128 & 1168.4 & 63.6 & 0.0755 \\
(1000,1) & A & $1.314\times10^{8}$ & 128--128 & 1104.0 & 128.0 & 0.0167 \\
(1000,1) & B & $1.311\times10^{8}$ & 128--128 & 1104.0 & 128.0 & 0.0132 \\
(1000,1) & C & $1.238\times10^{8}$ & 128--128 & 1104.0 & 128.0 & 0.0773 \\
(1000,1) & D & $1.183\times10^{8}$ & 128--128 & 1168.6 & 63.4 & 0.0837 \\
(1,1000) & A & $1.507\times10^{8}$ & 128--128 & 1104.0 & 128.0 & 0.0139 \\
(1,1000) & B & $1.516\times10^{8}$ & 128--128 & 1104.0 & 128.0 & 0.0142 \\
(1,1000) & C & $1.365\times10^{8}$ & 128--128 & 1104.0 & 128.0 & 0.0752 \\
(1,1000) & D & $1.348\times10^{8}$ & 128--128 & 1169.2 & 62.8 & 0.0759 \\
\bottomrule\end{tabular}
}
\end{table}

首先，矩阵条件数不包含真解到特征空间的距离。一个条件良好但空间表达能力不足的系统，也可能给出准确求解的错误近似。其次，相同行预算不保证相同的约束经验分布；矩阵列的可辨识性略有改善，可能伴随对界面约束覆盖的改变。再次，真解误差对应特定右端及预测位置，最坏方向未必被显著激发。

这与第二章固定矩阵误差分解一致：除了扰动放大项，还存在验证近似、截断和求值项。当前主比较未截断出明显秩亏，不能据此排除空间与经验目标的影响；后续增宽实验出现的有效秩不足则属于另一种情形。统一使用条件数解释所有实验，会混淆不同宽度和截断状态。

时间方面，自适应方案约为多次求解及候选评估之和，而A/B只求解一次。表中小规模CPU时间受系统状态影响，不宜将单次差值解释为精确硬件常数；但完整成本显著增加且解误差没有改善，足以限制当前方案的实用性主张。更严格的固定峰值串行重测已单独报告，不能拿并发阶段的时间去宣称加速。

\subsection{配额轨迹提供的机制证据}
表\ref{tab:history}列出$\beta=(1000,1)$时D的三轮平均轨迹，均值不必是整数。每次实际运行中的配额均为整数并满足预算，不能把均值中的小数理解为程序使用了小数个点。
\begin{table}[htbp]\centering
\caption{联合策略D的实际平均配额轨迹；每侧64特征、五种子}\label{tab:history}
\begin{tabular}{lllll}\toprule
预算 & 轮次 & 左内部均值 & 右内部均值 & 界面均值 \\ \midrule
点 & 0 & 128.0 & 128.0 & 32.0 \\
点 & 1 & 212.4 & 225.6 & 42.0 \\
点 & 2 & 307.0 & 310.2 & 54.8 \\
点 & 3 & 399.0 & 397.0 & 68.0 \\
行 & 0 & 128.0 & 128.0 & 32.0 \\
行 & 1 & 223.0 & 238.6 & 41.2 \\
行 & 2 & 331.2 & 332.8 & 52.0 \\
行 & 3 & 438.6 & 426.6 & 63.4 \\
\bottomrule\end{tabular}

\end{table}

点预算下最终界面平均68点，少于静态128点，而内部各侧平均接近400点；行预算下界面平均约63.4点。对应物理点均值超过1104，是因为一个界面位置占两行，减少界面后可用同样行预算放更多内部位置。这一计数关系解释了表\ref{tab:rankcost}中D的物理点更多，却没有使用额外标量行。

轨迹支持“当前指标倾向把资源给内部”这一描述，但不单独证明界面点少就是误差变大的唯一原因。改变配额同时改变位置分布、组经验平均和拟合结果；还可能存在某些内部点重复提供近似信息。后续同配额固定随机对照只回答一项有限问题，不能完成全部机制识别。

如果以后改进分配规则，首先应规定如何平衡物理通量与归一化通量，以及是否按每组单位成本的收益分配。当前按残差幅值分配并未测量“再增加一点能使真解误差降低多少”，两者不是同一量。由于真实应用没有真解，任何收益估计也必须来自可计算的代理，而不能把本实验已知解析解偷偷用于选点。

\section{参数敏感性结果的进一步解释}
宽度、权重范围和点数不是可以独立保证收敛的三个旋钮。增宽扩展名义特征空间，但若新增列高度相关，可辨识维数增长有限；增加点数可能改善覆盖，却因组平均权重改变和空间误差重新分配而导致解误差非单调；改变范围既改变特征尺度，也改变导数和矩阵谱。因此单因素扫描回答的是固定其余设置时的局部响应，不是完整的联合最优化。

在平滑直线模型上较小范围表现较好，可以解释为某些低频结构容易由缓变特征组合表示；同时小范围也可能增强列相关。这两种现象并不冲突。若只报告误差，可能忽略系数病态；若只报告条件数，又可能错过一个在指定双精度环境下确实准确的近似。本文把二者并列，避免用单个数字替代数值分析。

界面平方权重增加可加强满足界面约束的倾向，但可能牺牲内部拟合。它与增加界面点数并不等价：前者直接改变组目标的重要性，后者在组均方归一化下主要改变经验覆盖。把二者都笼统称为“加强界面损失”，会掩盖具体机制。本文所有对照均保留原权重与点数，便于单独核对。

岭参数也不是越大越稳定越好。较强收缩可以减小系数，但可能去掉真解所需方向；未经截断和经过岭的条件数必须区分数据矩阵与增广矩阵。现有扫描没有建立自动选择岭参数的方法，也没有用最终解析误差为每个种子挑参数，因此报告的是固定设置的敏感性，而非可在未知真解下直接部署的参数选择规则。

\section{近似复现中的可重复与可对齐}
可重复是指用本项目公开配置能够再次得到相近输出；与原文可对齐则要求使用相同模型、初始化、点分布、导数、权重、求解阈值和误差规则。前者并不自动推出后者。本文已经保存自身配置，但缺少作者若干实现细节，因此保留近似复现定位。

原表误差与本项目差异不能通过任意更换种子、删去较差运行或只保留较小积分值来消除。先前对初始化范围的另一解释也保留为诊断，而没有改写成作者原始方案。这种处理可能使结果表更复杂，却能使读者知道哪些设置来自原文字面，哪些来自本项目的补充假设。

解析导数与FD导数在部分设置接近，只说明在当前初始化、阈值及误差量级下，两种实现得到相近预测；不能证明FD矩阵与解析矩阵完全相同，也不能推导所有范围和步长下等价。步长扫描显示很小的差分步长可能恶化误差，符合舍入放大的机制，但其定量幅度还受线性系统敏感性影响。

花形积分误差的相对变化在最小误差组较大，提示接近极小数值时应降低对末几位的解释强度。本文不比较两方法某一末位数字来宣称优劣，也不把数值积分结果当作严格误差认证。后续若需高精度结论，应增加积分密度和高精度求值，并明确把额外成本计入。

\section{对原研究设想的回答}
原设想是利用内部与界面残差共同分配资源，以便在界面困难或高对比度时更有效地采样。现有实现确实完成了这一机制，也完成了相应基线及消融；因此研究过程仍围绕原问题展开。变化在于结果没有支持预期的稳定优势，论文据此把贡献从候选新方法收束为独立实现与系统分析。

已有证据支持三个有限判断：一是合理静态配额有时已能达到消融收益，二是内部残差下降不保证解误差和物理通量同步下降，三是自适应本身有不可忽略的完整成本。这些判断都能在表格中找到对应配置，不依赖作者主观挑选最好运行。它们也不等价于证明更精细的联合指标永远无用。

后续最有价值的改进应针对已观察的机制提出可证伪假设，例如在固定表达空间下分离配额与位置选择，或在冻结权重后使用新的监测池。若改成增长网络、更多子域或更大特征数，必须给静态基线相同新增表达能力，否则无法知道收益是否来自采样。先明确比较对象，再决定是否值得继续扩展，比不断扩大参数扫描更有解释力。

本章最终没有以“算法一定成功”作为验收标准，而是以模型与代码一致、资源计数清楚、全种子数据保留和结论可核对作为标准。尚未满足的学校形式要求及文献缺口属于终稿准备事项，不能通过把负面结果改成正面陈述来解决。

